\chapter{可微体渲染的地图表示}

\textbf{作者：} Hidenobu Matsuki 和 Andrew J. Davison

底层场景表示的选择会显著影响 SLAM 系统的能力和效率。如第 5 章所述，SLAM 中每种场景或地图表示都各有优缺点，具体取决于相机自运动跟踪、稠密场景重建或对象级分解等下游任务。近年来，得益于图形硬件和自动微分框架的进步，可通过可微渲染优化的三维场景表示，例如 neural radiance field（NeRF）\textsuperscript{[764]} 和 3D Gaussian Splatting（3DGS）\textsuperscript{[552]}，已成为计算机图形学界进行高质量逆渲染和新视图合成的重要选择。

这些表示最初面向已知相机位姿条件下的离线新视图合成，后来逐渐被用于三维重建和场景理解的多种其他任务 \textsuperscript{[1202]}，包括大规模场景的准确稠密几何重建 \textsuperscript{[819, 1241, 1159]}、语义融合 \textsuperscript{[1291, 1292, 586]} 等。早期实现需要数天优化；通过利用体素或点云等显式数据结构实现快速数据访问，优化时间已降至接近实时。近来，3D Gaussian Splatting 提供了极快且高质量的渲染能力，已经在许多应用中有效替代 Neural Radiance Fields。这些新的图形基元也正在被纳入 SLAM 研究。尽管将可微渲染与三维占据地图结合的思想可追溯至 2000 年代初 Thrun 等人提出的“前向传感器模型” \textsuperscript{[1089]}，现代硬件和算法正在开启全新的可能性。

本章概述 NeRF 和 3DGS 作为 SLAM 场景表示的方法，讨论其优缺点并提出未来研究方向。我们先回顾三维场景表示的历史（第 14.1 节），再详细介绍 NeRF（第 14.2 节）和 3DGS（第 14.3 节），包括它们与 SLAM 的集成，最后讨论潜在未来方向（第 14.4 节）。

\section{三维场景表示与可微渲染}

\subsection{可学习三维表示}

随着深度学习在二维图像领域取得成功 \textsuperscript{[260, 443, 444]}，研究界逐步探索用于三维任务的学习方法。早期工作直接使用点 \textsuperscript{[890]}、占据栅格 \textsuperscript{[929]} 等经典数据结构表示的三维数据，并将其用于下游任务。然而，这些经典表示是离散的，通常不适合深度学习：它们难以以紧凑方式表示拓扑灵活的水密连续表面，从而限制了神经网络训练所必需的基于梯度的优化。为构建更利于优化的三维表示，一些早期工作使用神经网络作为标量函数近似器来定义占据函数或有符号距离函数 \textsuperscript{[759, 847, 192]}，并通过直接三维监督训练。举例来说，DeepSDF \textsuperscript{[847]} 如图 14.1 所示，采用八层多层感知机（MLP），直接从三维坐标和可选潜变量回归有符号距离函数（SDF）。这些方法展示了不受拓扑和离散化误差约束的连续、紧凑三维形状重建。但三维数据的监督训练需要昂贵的真实标注，难以扩展数据集采集。这一限制促使研究者利用更易获得的二维图像信息和不同程度的监督进行三维场景理解，进而发展出通过可微渲染实现的三维重建。

\begin{figure}[htbp]
\centering
\includegraphics[width=\linewidth,height=0.78\textheight,keepaspectratio]{535f63a08501990b27c12ae24be5011e6ca98427910d276d7cac6dd4eb5c8a09.jpg}
\label{fig:chapter-14-535f63a08501990b27c12ae24be5011e6ca98427910d276d7cac6dd4eb5c8a09}
\end{figure}

\begin{figure}[htbp]
\centering
\includegraphics[width=\linewidth,height=0.78\textheight,keepaspectratio]{11d4928b8b2e1c332b8514e323e2d99c4934ba59f74308e3a244ce107b8a609d.jpg}
\label{fig:chapter-14-11d4928b8b2e1c332b8514e323e2d99c4934ba59f74308e3a244ce107b8a609d}
\end{figure}

\begin{figure}[htbp]
\centering
\includegraphics[width=\linewidth,height=0.78\textheight,keepaspectratio]{d34a5e7f1f88eb2ca16c0bab1ca84d31e686917ec36bed663ade65f14316429c.jpg}
\label{fig:chapter-14-d34a5e7f1f88eb2ca16c0bab1ca84d31e686917ec36bed663ade65f14316429c}
\end{figure}

(a) 体素
\begin{figure}[htbp]
\centering
\includegraphics[width=\linewidth,height=0.78\textheight,keepaspectratio]{b7b65df713b25cb7d451e855b169404e3ecad68a5699d48ffb26ebaf9bf4807c.jpg}
\label{fig:chapter-14-b7b65df713b25cb7d451e855b169404e3ecad68a5699d48ffb26ebaf9bf4807c}
\end{figure}

\begin{figure}[htbp]
\centering
\includegraphics[width=\linewidth,height=0.78\textheight,keepaspectratio]{22453fec121dc1dd2efc4dee6a34b47dac03d307382330695212f7b763c84955.jpg}
\label{fig:chapter-14-22453fec121dc1dd2efc4dee6a34b47dac03d307382330695212f7b763c84955}
\end{figure}

(b) 点
\begin{figure}[htbp]
\centering
\includegraphics[width=\linewidth,height=0.78\textheight,keepaspectratio]{f6fa85411af9d70f1c739dea526b7ea225307e1282da7196a0b465c40b18d9ea.jpg}
\label{fig:chapter-14-f6fa85411af9d70f1c739dea526b7ea225307e1282da7196a0b465c40b18d9ea}
\end{figure}

(c) 网格

\begin{figure}[htbp]
\centering
\includegraphics[width=\linewidth,height=0.78\textheight,keepaspectratio]{b269d6987a7ad91db25bdd1836cfc0a28c5ebb468cdd90237aef4900bfb4b728.jpg}
\label{fig:chapter-14-b269d6987a7ad91db25bdd1836cfc0a28c5ebb468cdd90237aef4900bfb4b728}
\end{figure}

\begin{figure}[htbp]
\centering
\includegraphics[width=\linewidth,height=0.78\textheight,keepaspectratio]{cee8e0cb4032ea415b61772648306ebc6b72f983f2c6b169226acb2ce4d19a5e.jpg}
\label{fig:chapter-14-cee8e0cb4032ea415b61772648306ebc6b72f983f2c6b169226acb2ce4d19a5e}
\end{figure}

(a) Occupancy Networks \textsuperscript{[759]}
(d) 本文方法

\begin{figure}[htbp]
\centering
\includegraphics[width=\linewidth,height=0.78\textheight,keepaspectratio]{4a992c6c1d969ceb41711d594dfda8d0694308faa1616bc181e330c7cec2fdf0.jpg}
\label{fig:chapter-14-4a992c6c1d969ceb41711d594dfda8d0694308faa1616bc181e330c7cec2fdf0}
\end{figure}

\begin{figure}[htbp]
\centering
\includegraphics[width=\linewidth,height=0.78\textheight,keepaspectratio]{f7c6686e4decf4c9e720c9e7b406004a5594478bd5287705ce8c9005a74fea99.jpg}
\label{fig:chapter-14-f7c6686e4decf4c9e720c9e7b406004a5594478bd5287705ce8c9005a74fea99}
\end{figure}

\begin{figure}[htbp]
\centering
\includegraphics[width=\linewidth,height=0.78\textheight,keepaspectratio]{2fb3de799b73d73abcabfec8ab7adcdefd1e864c77e8dbf0571f336ee33e35b4.jpg}
\label{fig:chapter-14-2fb3de799b73d73abcabfec8ab7adcdefd1e864c77e8dbf0571f336ee33e35b4}
\end{figure}

(b) DeepSDF \textsuperscript{[847]}
图 14.1 可学习三维表示。ONet \textsuperscript{[759]}（©2019 IEEE）和 DeepSDF \textsuperscript{[847]}（©2019 IEEE）都使用神经网络回归连续隐式表面函数。

\subsection{可微渲染}

三维数据的监督学习面临大规模真实标注难以获取且成本高昂的挑战。为克服这些问题，研究者将图形渲染流水线纳入神经网络训练，从而推动了可微渲染（differentiable rendering，DR）的发展。DR 允许梯度从二维图像反向传播到三维场景参数（图 14.2）。在图形流水线中，三维模型可以采用两种表示：

- \textbf{表面数据}：显式编码对象边界，例如多边形或隐式表面。
- \textbf{体积数据}：描述密度或颜色等连续场，表示对象的内部属性和材料分布。

这两类表示对应的渲染过程分别称为表面渲染和体渲染。因此，各自采用不同的渲染方法（图 14.3）：

- \textbf{表面渲染}通常采用光栅化，将表面几何投影到二维平面。它可利用硬件加速，因此对显式表面非常高效。
- \textbf{体渲染}通常通过光线步进实现：在体积中投射光线，并沿每条光线路径采样数值。它适于连续场，但计算代价较高。

尽管早期 DR 方法对网格使用软化光栅化 \textsuperscript{[695, 541, 680]}，体渲染已逐渐成为更强大的基于梯度优化方法。体渲染在整个三维空间提供梯度，甚至在确定精确表面位置之前也是如此，从而使神经网络训练更容易，因此被广泛作为 DR 的骨干。继 NeuralVolume \textsuperscript{[688]} 之后，可微体渲染 \textsuperscript{[807]} 和 Scene Representation Networks \textsuperscript{[1015]} 等工作使用可微光线步进表示隐式表面。尽管最初仅限于简单形状，这些方法为更复杂、更真实的三维重建奠定了基础。

\subsection{使用可微渲染的三维场景表示}

2020 年，Mildenhall 等人 \textsuperscript{[764]} 提出 NeRF，成为一个重要突破。该方法从带位姿的彩色图像出发，使用体积光线步进优化 MLP 以重建场景。NeRF 使用单个 MLP 作为场景表示，沿像素光线步进并查询 MLP 的不透明度和颜色来执行体渲染。由于体渲染天然可微，MLP 能够利用多视图信息最小化渲染损失，从而实现高质量新视图合成。NeRF 的 MLP 输入包括空间位置和视角方向，并配合位置编码捕获场景高频细节。虽然 NeRF 获得了高质量新视图合成，但训练速度是其主要瓶颈。为此，后续工作引入体素栅格 \textsuperscript{[1050, 341]} 或哈希函数 \textsuperscript{[784]} 等显式体积结构来提高性能。这些进展表明，高质量新视图合成的关键在于可微体渲染，而非神经网络本身；不依赖 MLP 也可以达到与 NeRF 相当的渲染质量。不过，即使在这些改进系统中，逐像素光线步进仍是渲染速度的重要瓶颈。

\begin{figure}[htbp]
\centering
\includegraphics[width=\linewidth,height=0.78\textheight,keepaspectratio]{2998a1a838850aa3a70f4a7fc5231f4df5f76a655645b365470b36400adc0794.jpg}
\caption{图 14.2 可微渲染流水线概览。由于梯度可一致传播，DR 能够以端到端方式从多视图二维图像优化三维场景参数。}
\label{fig:chapter-14-2998a1a838850aa3a70f4a7fc5231f4df5f76a655645b365470b36400adc0794}
\end{figure}

\begin{figure}[htbp]
\centering
\includegraphics[width=\linewidth,height=0.78\textheight,keepaspectratio]{59e2786f7b5f5e1ed9962e22dd23d7e7cacd1986cc3a7450ba3647565a89eb16.jpg}
\label{fig:chapter-14-59e2786f7b5f5e1ed9962e22dd23d7e7cacd1986cc3a7450ba3647565a89eb16}
\end{figure}

(a) 光线步进（NeRF）

\begin{figure}[htbp]
\centering
\includegraphics[width=\linewidth,height=0.78\textheight,keepaspectratio]{71ed37f3db1397088d2db551b6cae7ba28bd50c8d950e90faed43b05208205f4.jpg}
\label{fig:chapter-14-71ed37f3db1397088d2db551b6cae7ba28bd50c8d950e90faed43b05208205f4}
\end{figure}

(b) 光栅化/Splatting（3DGS）
图 14.3 NeRF 和 3D Gaussian Splatting 渲染方法的比较。基于光线步进的方法混合沿射线固定间隔采样点的推理结果；光栅化则将射线上的每个图元投影到图像平面后再混合。

在 NeRF 和体渲染快速发展的同时，研究者也探索了可微的基于点渲染技术 \textsuperscript{[597, 954, 556]}。这些方法采用点云 \textsuperscript{[597, 954]} 或 metaball \textsuperscript{[556]} 等基于点的表示，并通过基于 CNN 的光栅化器或近似渲染器定义可微渲染流水线。2023 年，Zwicker 等人 \textsuperscript{[1319]} 提出了一种基于 3D Gaussian Splatting \textsuperscript{[552]}（3DGS）的截然不同的点渲染方法，用于辐射场捕获。3DGS 使用大量半透明高斯团表示场景，并通过 alpha 混合渲染。不同于 NeRF，3DGS 使用可微光栅化（图 14.3）。它和传统图形光栅化一样，遍历待光栅化的图元而非沿射线步进，因而利用三维场景的天然稀疏性，既能表达高保真三维场景，又能在保持与最高效体积 NeRF 相近训练速度的同时显著加快渲染。

下文将探讨 NeRF 和 3DGS 的技术细节（图 14.4）。尽管二者共同目标是辐射场捕获，但场景表示和渲染流水线存在显著差异。我们先概述 NeRF，重点介绍其基于光线步进的可微渲染表述及其使用不同数据结构、发展为通用 neural field 表示的扩展；随后考察 3DGS 及其基于光栅化的渲染流水线，并强调其与 NeRF 的关键区别。

\begin{figure}[htbp]
\centering
\includegraphics[width=\linewidth,height=0.78\textheight,keepaspectratio]{37a1c71b670f48f773a63b744b5e791df04edb75800d6593b48265332a5e6d5a.jpg}
\label{fig:chapter-14-37a1c71b670f48f773a63b744b5e791df04edb75800d6593b48265332a5e6d5a}
\end{figure}

\begin{figure}[htbp]
\centering
\includegraphics[width=\linewidth,height=0.78\textheight,keepaspectratio]{0f8316492b7b4d357f4ab324d007d57241ba82bfc02b333d295e65a792b83d9c.jpg}
\caption{图 14.4 MipNeRF360 \textsuperscript{[58]}（左，©2019 IEEE）和 3DGS \textsuperscript{[552]}（右）的新视图合成示例。两种方法都可从真实世界中带位姿的 RGB 图像捕获高保真辐射场。图像来自 \textsuperscript{[58, 552]}。}
\label{fig:chapter-14-0f8316492b7b4d357f4ab324d007d57241ba82bfc02b333d295e65a792b83d9c}
\end{figure}

\section{Neural Radiance Fields（NeRF）}

本节说明 neural radiance field（NeRF）的技术细节、衍生方法及其在视觉 SLAM 中的应用。这些技术背景为理解下一节中与 3D Gaussian Splatting 的动机和差异提供基础。

\subsection{方法概述}

\begin{figure}[htbp]
\centering
\includegraphics[width=\linewidth,height=0.78\textheight,keepaspectratio]{eaf32196cd7b30e49794569a6ee38afdaf6585d7212cd6c6647c20006ce3fac9.jpg}
\caption{图 14.5 NeRF 优化流水线概览。辐射场由 MLP 表示，通过可微体渲染生成二维图像。优化 MLP 参数以最小化渲染彩色图像与真实彩色图像之间的光度误差。图像来自 \textsuperscript{[764]}（©2020 Springer）。}
\label{fig:chapter-14-eaf32196cd7b30e49794569a6ee38afdaf6585d7212cd6c6647c20006ce3fac9}
\end{figure}

NeRF 模型表示为函数 $f_\Theta(\mathbf x,\mathbf d)\to(\mathbf c,\sigma)$，其中 $\Theta$ 是 MLP 权重，场景中的三维坐标为 $\mathbf x=(x,y,z)$，视角方向为 $\mathbf d=(\theta,\phi)$，颜色为 $\mathbf c=(r,g,b)$，体积密度为 $\sigma$。为合成新视图，NeRF 向场景投射相机光线，为每个像素生成采样点。使用 NeRF MLP 在每个采样点计算局部颜色和密度，再通过体渲染合成二维图像（图 14.5）。相机光线 $\mathbf r(t)=\mathbf o+t\mathbf d$ 的期望颜色 $C(\mathbf r)$ 按下列积分计算：

\begin{equation}
C (\mathbf {r}) = \int_ {t _ {n}} ^ {t _ {f}} T (t) \sigma (\mathbf {r} (t)) \mathbf {c} (\mathbf {r} (t), \mathbf {d}) \mathrm{d} t.\tag{14.1}
\end{equation}

其中 $dt$ 表示每个积分步中光线行进的微分距离，$t_n$ 和 $t_f$ 分别表示近、远边界。$\sigma(\mathbf r(t))$ 和 $\mathbf c(\mathbf r(t),\mathbf d)$ 分别表示具有视角方向 $d$ 的相机光线上点 $\mathbf r(t)$ 的体积密度和颜色。从 $t_n$ 到 $t$ 的累积透射率为 $T(t)=\exp(-\int_{t_n}^t\sigma(\mathbf r(s))\,\mathrm{d}s)$。连续积分通过数值求积进行估计：将光线划分为 $N$ 个等间隔区间，在区间之间均匀采样并在每个区间内随机采样。利用求积规则，体渲染方程写作：

\begin{equation}
\hat {C} (\mathbf {r}) = \sum_ {i = 1} ^ {N} \alpha_ {i} T _ {i} \mathbf {c} _ {i}, \quad \text { where } T _ {i} = \exp \Bigl (- \sum_ {j = 1} ^ {i - 1} \sigma_ {j} \delta_ {j} \Bigr).\tag{14.2}
\end{equation}

$\delta_i=t_{i+1}-t_i$ 表示连续采样点 $t_i$ 与 $t_{i+1}$ 之间的间隔，$\sigma_i$ 和 $\mathbf c_i$ 分别表示在光线上第 $i$ 个采样点求得的密度和颜色；此外，$\alpha_i=1-\exp(-\sigma_i\delta_i)$ 表示第 $i$ 个采样点通过 alpha 合成得到的不透明度。

为优化模型，使用渲染图像与真实彩色图像之间的平方光度误差：

\begin{equation}
L = \sum_ {r \in R} \| \hat {C} (\mathbf {r}) - C _ {g t} (\mathbf {r}) \| _ {2} ^ {2},\tag{14.3}
\end{equation}

其中 $C_{gt}(r)$ 表示训练图像中与 $r$ 关联像素的真实颜色，$R$ 表示用于合成目标图像的一批光线。在 SLAM 中，常利用不透明度的 alpha 混合渲染深度，从而额外最小化期望深度与传感器深度或预测深度之间的差异。

\subsection{NeRF 中的数据结构}

原始 NeRF 论文 \textsuperscript{[764]} 使用单个 MLP 表示整个辐射场；在特定三维位置评估颜色和密度时，需要一次前向推理。这意味着所有网络参数都参与评估，导致前向推理和反向训练都非常慢。为加快推理和训练，多个方法部分或完全引入显式数据结构，实现快速数据访问。

\textbf{体素/哈希栅格。} 早期工作引入体素栅格以加速前向推理。随后，Plenoxel \textsuperscript{[341]} 和 DirectVoxGo \textsuperscript{[1050]} 将体素特征栅格作为可训练属性，并将辐射场编码其中。与基于 MLP 的 NeRF 相比，特定位置的颜色和密度只需简单索引查询和三线性插值，显著加速整个流水线（从数天降至 10 分钟）。Instant-NGP \textsuperscript{[784]}（图 14.6）进一步引入多分辨率哈希编码，在哈希表中高效编码体积特征；结合高度优化的 CUDA 实现，可在数秒内完成 NeRF 训练。

\textbf{点和 surfel。} 基于点的表示也可与 NeRF 结合使用（图 14.7）。Point-NeRF \textsuperscript{[1209]} 优化带特征的点云和微型 MLP 解码器，并通过体渲染框架渲染图像。它表明，相比规则体素栅格，点表示具有空间分配灵活性，能够跳过空白区域，从而比原始 NeRF 快三倍。类似点云，也可以用 surfel 表示辐射场 \textsuperscript{[366]}。

\begin{figure}[htbp]
\centering
\includegraphics[width=\linewidth,height=0.78\textheight,keepaspectratio]{fb41999f11dc2d80bb595dd8d3702ef7dae6b700f21955a48554097487e8153e.jpg}
\caption{图 14.6 基于栅格的 NeRF 示例（Instant-NGP \textsuperscript{[784]}）。辐射场由多分辨率哈希栅格和微型 MLP 解码器的组合建模。图像来自 \textsuperscript{[784]}。}
\label{fig:chapter-14-fb41999f11dc2d80bb595dd8d3702ef7dae6b700f21955a48554097487e8153e}
\end{figure}

\begin{figure}[htbp]
\centering
\includegraphics[width=\linewidth,height=0.78\textheight,keepaspectratio]{3c4b58f41b467e17a22c910d82f1a60a10fdf839ca4fd3d89902ad79e661ae33.jpg}
\caption{图 14.7 基于点的 NeRF 表示示例（Point-NeRF \textsuperscript{[1209]}）。辐射场由带可训练特征的点与微型 MLP 解码器的组合建模。图像来自 \textsuperscript{[1209]}（©2022 IEEE）。}
\label{fig:chapter-14-3c4b58f41b467e17a22c910d82f1a60a10fdf839ca4fd3d89902ad79e661ae33}
\end{figure}

\subsection{Neural Fields}

早期 NeRF 工作主要关注将三维位置映射到颜色和密度的辐射场，以实现照片级真实感场景捕获；但根据下游任务不同，这一基本思想也可扩展到各种其他输出。这类通用场表示称为 Neural Fields。下面介绍几个主要例子。

\textbf{表面表示（SDF 与占据场）。} 为从 RGB 图像获得几何上更准确的三维重建，Neural Fields 可以使用有符号距离函数（SDF）和占据场等表示（图 14.8）。NeuS \textsuperscript{[1159]} 和 VolSDF \textsuperscript{[1241]} 等技术利用 SDF 表示表面。SDF 编码点到最近表面的距离，符号表示点位于对象内部还是外部。将 SDF 集成到 NeRF 中可以实现精确的表面建模。UNISURF \textsuperscript{[819]} 等技术使用占据场表示三维空间中表面的存在与否：占据场将空间划分为栅格，每个单元表示是否被表面占据。该方法可以与 NeRF 结合，有效处理复杂场景。为从这些表示中提取网格，可以使用 marching cubes 等技术，将连续表面表示转换为离散网格格式。

\begin{figure}[htbp]
\centering
\includegraphics[width=\linewidth,height=0.78\textheight,keepaspectratio]{b2db701fdc19aa6fea11df4f1d15f4a1c0755bdee861bf8c7636e420fb837ef4.jpg}
\caption{图 14.8 用于表面重建的 Neural Field。UNISURF 使用占据场，NeuS 使用由 MLP 建模的连续有符号距离函数；二者都可通过体渲染生成图像。图像来自 \textsuperscript{[819, 1159]}。}
\label{fig:chapter-14-b2db701fdc19aa6fea11df4f1d15f4a1c0755bdee861bf8c7636e420fb837ef4}
\end{figure}

\textbf{语义。} 将语义信息纳入 Neural Fields，可以把场景分解为有意义的组件，增强场景理解和交互能力（图 14.9）。早期工作 \textsuperscript{[1291, 1292]} 扩展 NeRF，使其除颜色和密度外还输出语义标签：利用多视图语义图像训练模型，使其能够预测场景每一部分的语义类别。这些场将多视图的高维特征融合为三维一致的语义特征场。更一般地，NeRF 可以融合预训练基础网络的二维特征，通过语言提示或点击执行场景分解和交互式编辑等更高级的语义任务 \textsuperscript{[586, 1111, 745]}。

\subsection{用于视觉 SLAM 的 NeRF/Neural Fields}

Neural Fields 最初设计用于利用带位姿的彩色图像进行离线新视图合成。然而，它们生成高保真重建的能力引发了将其用于视觉 SLAM 的兴趣，后者需要实时联合估计相机位姿和场景表示。

与经典稠密 SLAM 方法相比，基于 Neural Field 的三维表示具有多项优势。第一，通过可微渲染进行端到端优化，简化了二维到三维的数据融合，无需像基于 TSDF 的 SLAM 那样进行多阶段手工设计的融合步骤，而这些步骤的一致性并不总能保证。第二，Neural Fields 可以建模复杂光场，更灵活地处理透明物体和反射，从而最小化光度误差。第三，根据数据结构不同，它们能够借助神经表示中固有的先验完成场景补全和压缩。

\begin{figure}[htbp]
\centering
\includegraphics[width=\linewidth,height=0.78\textheight,keepaspectratio]{0fdbfd192edb240482e4ae41ac28c34a11a9030bed44e549ba6a5f03d3cdac5f.jpg}
\caption{图 14.9 用于语义场景重建的 Neural Field。通过为 MLP 增加语义头，NeRF 可以输出多视图一致的语义标签/特征，并通过逐视图语义预测进行训练。图像来自 \textsuperscript{[1291, 586]}。}
\label{fig:chapter-14-0fdbfd192edb240482e4ae41ac28c34a11a9030bed44e549ba6a5f03d3cdac5f}
\end{figure}

iMAP \textsuperscript{[1047]}（图 14.10）是这一方向的首项工作，使用类似原始 NeRF 的纯 MLP 地图表示。iMAP 的关键洞见是，得益于 MLP 固有的平滑先验，基于 MLP 的方法可以实现地图压缩和场景补全。与 NeRF 研究的发展类似，研究者也在 SLAM 中探索经典表示和神经表示的结合。NICE-SLAM \textsuperscript{[1309]} 首次引入混合表示，将多分辨率体素特征栅格与微型 MLP 结合。其他工作也采用了类似的栅格表示 \textsuperscript{[1232, 525, 1153]}。NeRF-SLAM \textsuperscript{[940]} 利用 DROID-SLAM 的深度信息（参见第 13 章）加速 NeRF 重建。最近，Point-NeRF 风格的点基元也被应用于 SLAM \textsuperscript{[968]}。在诸多技术组件中，底层数据结构的选择对 SLAM 的重建质量、计算效率和实时性能影响最大。因此，下面考察 Neural-Field SLAM 方法中使用的各种数据结构，并详细说明其主要特征与权衡。

\subsubsection{Neural Field SLAM 中数据结构的特征与权衡}

底层数据结构的选择会显著影响基于 Neural Field 的 SLAM 属性。表 14.1 概览了各种 Neural-Field SLAM 方法。下面详细介绍这些表示相关的主要特征和权衡。

\textbf{压缩与推理速度之间的权衡。} 早期 Neural Field SLAM 使用单个 MLP，较新的方法则部分或完全使用类似经典图形基元的显式数据结构。MLP 用网络参数近似整个辐射场函数，在压缩方面具有巨大优势，但参数的物理含义难以解释。另一方面，体素栅格等显式数据结构更易解释，能够快速访问数据并局部更新地图，但需要更多内存。因此，地图表示体现了典型的时间--空间复杂度权衡。对于 SLAM 任务，实时运行尤其重要，实践中通常更偏好显式数据结构。

\begin{figure}[htbp]
\centering
\includegraphics[width=\linewidth,height=0.78\textheight,keepaspectratio]{c662a7a1daeb4a4439762f79c4c27d1961580b5c16b40cf367dc07d8dc330073.jpg}
\caption{图 14.10 iMAP \textsuperscript{[1047]} 的定性结果。该方法首次展示了可以使用 MLP 和可微渲染框架，从 RGB-D 视频流实时联合估计相机位姿和场景几何。}
\label{fig:chapter-14-c662a7a1daeb4a4439762f79c4c27d1961580b5c16b40cf367dc07d8dc330073}
\end{figure}

\small
\setlength{\tabcolsep}{2pt}
\begin{center}
\begin{tabular}{p{0.237\linewidth}p{0.237\linewidth}p{0.237\linewidth}p{0.237\linewidth}}
\toprule
底层数据结构 & MLP & 体素/哈希 & 点/surfel \\
\midrule
压缩 & 好 & 中等 & 中等 \\
推理速度 & 受限 & 好 & 中等 \\
场景补全 & 好 & 好 & 中等 \\
动态分配 & 受限 & 受限 & 好 \\
遗忘问题 & 受限 & 中等 & 中等 \\
\bottomrule
\end{tabular}
\end{center}

表 14.1 Neural Field 的底层数据结构比较。

\textbf{场景补全。} iMAP 表明，由于函数的连续性，基于 MLP 的三维重建能够填补未观测区域的孔洞。多分辨率栅格也具有类似属性（图 14.11），因为位置编码或多分辨率栅格会在频率或栅格分辨率内平滑信号。基于点的表示则没有这种效果，因为其基元是离散的，不具备全局平滑能力。

\textbf{图形基元的动态分配。} 与离线场景重建任务不同，SLAM 需要根据运行时观测在线更新状态变量，以适应场景边界变化和回环等动态变化。在 Neural Fields 中，MLP 和体素类型表示由于边界和位置编码/体素分辨率固定，难以进行这些更新，通常需要动态分配多个局部 NeRF 模型 \textsuperscript{[1070]}（图 14.12）。相比之下，点表示在基元分配方面更灵活，能够更有效地处理回环 \textsuperscript{[673]}。

\begin{figure}[htbp]
\centering
\includegraphics[width=\linewidth,height=0.78\textheight,keepaspectratio]{ed4ad87271b3dbdb9399c944314bf8eaf17062e0acc8d4b0a89f4b935c719cef.jpg}
\label{fig:chapter-14-ed4ad87271b3dbdb9399c944314bf8eaf17062e0acc8d4b0a89f4b935c719cef}
\end{figure}

\begin{figure}[htbp]
\centering
\includegraphics[width=\linewidth,height=0.78\textheight,keepaspectratio]{7599c1debbca48b0237832506043f93f1531efbe99134fdf0e544950d2390341.jpg}
\label{fig:chapter-14-7599c1debbca48b0237832506043f93f1531efbe99134fdf0e544950d2390341}
\end{figure}

(a) iMAP \textsuperscript{[1047]}

\begin{figure}[htbp]
\centering
\includegraphics[width=\linewidth,height=0.78\textheight,keepaspectratio]{242b53d23bc51e8f01fc52aade3eca982bdaf21b3caa351056af13a25f2bdb41.jpg}
\label{fig:chapter-14-242b53d23bc51e8f01fc52aade3eca982bdaf21b3caa351056af13a25f2bdb41}
\end{figure}

\begin{figure}[htbp]
\centering
\includegraphics[width=\linewidth,height=0.78\textheight,keepaspectratio]{c3ce9f9460c49ce404c5f7d8ed23280e65accfc449c52b697b2e2c9e5c4eae5c.jpg}
\label{fig:chapter-14-c3ce9f9460c49ce404c5f7d8ed23280e65accfc449c52b697b2e2c9e5c4eae5c}
\end{figure}

(b) NICE-SLAM \textsuperscript{[1309]}
图 14.11 基于 NeRF 的 SLAM 的孔洞填充效果。得益于多分辨率频率或栅格编码以及连续 MLP 函数，该表示能够平滑局部几何，并为小范围未观测区域提供合理的孔洞填充效果。图像来自 \textsuperscript{[1047]} 和 \textsuperscript{[1309]}（©2022 IEEE）。

\begin{figure}[htbp]
\centering
\includegraphics[width=\linewidth,height=0.78\textheight,keepaspectratio]{53355a80ce2e503687bdb55968deeeac526d450ce56e1a5e0196158b3f29d460.jpg}
\label{fig:chapter-14-53355a80ce2e503687bdb55968deeeac526d450ce56e1a5e0196158b3f29d460}
\end{figure}

\begin{figure}[htbp]
\centering
\includegraphics[width=\linewidth,height=0.78\textheight,keepaspectratio]{f16aaf662901798164e25fbda400cf8167387f91f7f378fea85293595bc36c2c.jpg}
\label{fig:chapter-14-f16aaf662901798164e25fbda400cf8167387f91f7f378fea85293595bc36c2c}
\end{figure}

\begin{figure}[htbp]
\centering
\includegraphics[width=\linewidth,height=0.78\textheight,keepaspectratio]{14f610e64d11dca91850044136eaf345e4174360cbb7808d69d332fa577cfa92.jpg}
\label{fig:chapter-14-14f610e64d11dca91850044136eaf345e4174360cbb7808d69d332fa577cfa92}
\end{figure}

\begin{figure}[htbp]
\centering
\includegraphics[width=\linewidth,height=0.78\textheight,keepaspectratio]{c38c1f8c0696debadafdf79e5cfd147efaff19bc3e1f4aa6a33f02e570488e50.jpg}
\label{fig:chapter-14-c38c1f8c0696debadafdf79e5cfd147efaff19bc3e1f4aa6a33f02e570488e50}
\end{figure}

(a) Loopy-SLAM \textsuperscript{[673]}

\begin{figure}[htbp]
\centering
\includegraphics[width=\linewidth,height=0.78\textheight,keepaspectratio]{ea05e0f0951fa7bd80ebc5a259abd2652495a0b7c9de753ea4c3dd43f9a5152d.jpg}
\label{fig:chapter-14-ea05e0f0951fa7bd80ebc5a259abd2652495a0b7c9de753ea4c3dd43f9a5152d}
\end{figure}

\begin{figure}[htbp]
\centering
\includegraphics[width=\linewidth,height=0.78\textheight,keepaspectratio]{f1036b3f09c8948fd1249f57ca1a69878d9e148f7bffdc77520707488e9168a3.jpg}
\label{fig:chapter-14-f1036b3f09c8948fd1249f57ca1a69878d9e148f7bffdc77520707488e9168a3}
\end{figure}

\begin{figure}[htbp]
\centering
\includegraphics[width=\linewidth,height=0.78\textheight,keepaspectratio]{c285e4eafc541ee061910b9f3cec5d8dd2ae2052dae5ee7ba5109189d30da9fa.jpg}
\label{fig:chapter-14-c285e4eafc541ee061910b9f3cec5d8dd2ae2052dae5ee7ba5109189d30da9fa}
\end{figure}

(b) MIPS-Fusion \textsuperscript{[1070]}
图 14.12 基于 NeRF 的 SLAM 回环。基于体素的方法需要多个局部 Neural Field 子地图来处理回环导致的相机位姿更新，而离散的点表示可以更无缝地改变基元分配。图像来自 \textsuperscript{[673]}（©2024 IEEE）和 \textsuperscript{[1070]}（©2023 ACM）。

\textbf{遗忘问题。} 在纯 MLP 表示和混合表示中，全局函数的存在意味着任何局部参数更新都会影响整体结果。因此，地图必须在优化窗口中保留过去信息作为训练信号，并需要显式关键帧数据库。这个通常称为“遗忘问题”的现象，可以通过减少对全局函数（通常由 MLP 表示）的依赖，或完全移除全局函数来缓解。

\section{3D Gaussian Splatting}

\begin{figure}[htbp]
\centering
\includegraphics[width=\linewidth,height=0.78\textheight,keepaspectratio]{5a472fe82aec09b6085fdc15e33d573a0ff1a24732d855974e7e9596b43a141b.jpg}
\caption{图 14.13 3DGS 新视图渲染结果。（左）RGB 渲染；（右）高斯体着色。3DGS 通过分配大量半透明三维高斯体表示场景。}
\label{fig:chapter-14-5a472fe82aec09b6085fdc15e33d573a0ff1a24732d855974e7e9596b43a141b}
\end{figure}

NeRF 执行可微光线步进，而 3D Gaussian Splatting（3DGS）对体数据执行可微光栅化。类似常规图形光栅化，3DGS 遍历待光栅化的图元，而不是沿光线步进，利用三维场景的天然稀疏性，获得足以捕获高保真三维场景的表达能力，同时显著提高渲染速度（图 14.13）。输入是一组带有 SfM 估计位姿的彩色图像；此外，3DGS 使用 SfM 过程中生成的稀疏点云构造一组初始三维高斯图形基元。

\subsection{方法概述}

每个三维高斯由三维均值位置 $\pmb\mu\in\mathbb{R}^3$、协方差矩阵 $\pmb\Sigma\in\mathbb{R}^{3\times3}$、不透明度 $o$ 和由球谐（SH）表示的方向相关颜色 $c$ 参数化。给定缩放矩阵 $\pmb S\in\mathbb{R}^{3\times3}$ 和旋转矩阵 $R\in\mathrm{SO}(3)$，协方差矩阵 $\Sigma$ 表示为：

\begin{equation}
\boldsymbol {\Sigma} = \boldsymbol {R} \boldsymbol {S} \boldsymbol {S} ^ {\top} \boldsymbol {R} ^ {\top}.\tag{14.4}
\end{equation}

为实现可微光栅化，使用近似投影函数将这些三维高斯投影到图像平面（$z=1$ 平面）。投影高斯在二维图像空间中的协方差矩阵 $\Sigma'$ 为：

\begin{figure}[htbp]
\centering
\includegraphics[width=\linewidth,height=0.78\textheight,keepaspectratio]{906077fcec572428c124c7dc4f3f14f2efe5b6429251e18332c5b8071167b167.jpg}
\caption{图 14.14 3DGS 优化过程概览 \textsuperscript{[552]}。}
\label{fig:chapter-14-906077fcec572428c124c7dc4f3f14f2efe5b6429251e18332c5b8071167b167}
\end{figure}

\begin{equation}
\boldsymbol {\Sigma} ^ {\prime} = \mathbf {J} \mathbf {W} \boldsymbol {\Sigma} \mathbf {W} ^ {\top} \mathbf {J} ^ {\top},\tag{14.5}
\end{equation}

其中 $J$ 表示投影变换的雅可比矩阵，$W$ 是视图变换矩阵。将 $\Sigma'$ 左上角的 $2\times2$ 块作为二维协方差矩阵。通过查询像素 $\mathbf x_{2D}=[u,v]^\top$，可以得到二维高斯函数的值。

像素 $x_{2D}$ 的最终颜色 $C$ 由三维高斯的 alpha 混合确定：

\begin{equation}
C = \sum_ {i \in \mathcal {N}} \mathbf {c} _ {i} \alpha _ {i} \prod_ {j = 1} ^ {i - 1} (1 - \alpha _ {j}),\tag{14.6}
\end{equation}

其中 $\mathcal{N}$ 是沿像素光线的高斯集合，$\alpha_i$ 是不透明度 $o_i$ 与像素空间高斯函数的乘积：

\begin{equation}
\alpha_ {i} = o _ {i} G _ {\mathrm{2D}} \left(\mathbf {x} _ {2 D}\right).\tag{14.7}
\end{equation}

在 alpha 混合中，按照相对于渲染视点的深度顺序对三维高斯排序。输入图像被划分为 $16\times16$ 的图块，并使用 GPU 优化的基数排序对每个图块分别执行该排序。此外，为提高效率，只考虑位于当前视锥内的高斯（视锥剔除）。通过最小化渲染图像与真实图像之间 L1 损失和 SSIM 损失的加权和，估计三维高斯参数，权重为 $\lambda$：

\begin{equation}
\mathcal {L} = (1 - \lambda) \mathcal {L} _ {1} + \lambda \mathcal {L} _ {\mathrm{SSIM}}.\tag{14.8}
\end{equation}

此外，在优化过程中会根据高斯的梯度和不透明度执行自适应稠密化与剪枝，以灵活调整高斯数量（图 14.14）。

\subsection{3D Gaussian Splatting 的应用}

类似 NeRF 演变为通用 Neural Field，3DGS 也被应用于各种三维视觉任务。鉴于该领域研究数量庞大，这里重点介绍几篇早期工作。

\textbf{表面重建。} 3DGS 主要设计用于照片级真实感新视图合成，并不天然保证几何精度。若干工作针对这一局限进行了改进（图 14.15）。SuGAR \textsuperscript{[411]} 为 3DGS 引入几何正则化，实现几何准确的三维网格提取。该方法为每个三维高斯定义表面法向，并通过鼓励相邻高斯尽可能平坦来强制局部表面平滑。2D Gaussian Splatting（2DGS）\textsuperscript{[483]} 将三维高斯替换为二维圆盘，在图形基元中显式定义表面法向。通过加入法向和深度一致性损失，2DGS 进一步提高了表面重建精度。

\begin{figure}[htbp]
\centering
\includegraphics[width=\linewidth,height=0.78\textheight,keepaspectratio]{8b85e2de7c644747549b4d3ed93f8f88621313e2dd76a6d1b42e2df7b56db41d.jpg}
\label{fig:chapter-14-8b85e2de7c644747549b4d3ed93f8f88621313e2dd76a6d1b42e2df7b56db41d}
\end{figure}

(a) SuGAR \textsuperscript{[411]}（©2024 IEEE）

\begin{figure}[htbp]
\centering
\includegraphics[width=\linewidth,height=0.78\textheight,keepaspectratio]{de8a1cfcc5914cddc0cc5d4c83fc3dee49ebad32fefaff0e33f214cd13b46a1b.jpg}
\label{fig:chapter-14-de8a1cfcc5914cddc0cc5d4c83fc3dee49ebad32fefaff0e33f214cd13b46a1b}
\end{figure}

\begin{figure}[htbp]
\centering
\includegraphics[width=\linewidth,height=0.78\textheight,keepaspectratio]{3e57b8ac20d35d6452301a4c7cf04cd2249de5ab325cba4248f6f780a5965ec8.jpg}
\label{fig:chapter-14-3e57b8ac20d35d6452301a4c7cf04cd2249de5ab325cba4248f6f780a5965ec8}
\end{figure}

\begin{figure}[htbp]
\centering
\includegraphics[width=\linewidth,height=0.78\textheight,keepaspectratio]{4ca44d7e8516723b86a16ff3b549d8524d3ece77536a4fb69c902d69d56965ae.jpg}
\label{fig:chapter-14-4ca44d7e8516723b86a16ff3b549d8524d3ece77536a4fb69c902d69d56965ae}
\end{figure}

(b) 2DGS \textsuperscript{[483]}。

\begin{figure}[htbp]
\centering
\includegraphics[width=\linewidth,height=0.78\textheight,keepaspectratio]{7770a9b1c1367c02a5607c4caf847b55c7f98562e9b24d969ba86ed691af6e28.jpg}
\caption{图 14.15 面向表面的 Gaussian Splatting 方法示例。通过定义表面法向方向并强制相邻高斯之间保持平滑，这些方法能够实现准确的表面重建。图像来自 \textsuperscript{[411]}（©2024 IEEE）和 \textsuperscript{[483]}。}
\label{fig:chapter-14-7770a9b1c1367c02a5607c4caf847b55c7f98562e9b24d969ba86ed691af6e28}
\end{figure}

\textbf{语义场景理解。} 3DGS 也可以扩展到语义场景理解。与 NeRF 类似，最常见的做法是通过可微渲染支持的端到端训练，将二维语义分割传播到带语义特征的三维高斯中。例如，LangSplat \textsuperscript{[895]} 和 FMGS \textsuperscript{[1316]} 在 3DGS 训练过程中，将 CLIP、DINO 等二维基础模型集成到多视图一致的三维高斯中，实现开放词汇语义分割（图 14.16）。由于三维高斯是显式离散表示，编辑重建场景（例如操纵或删除）比连续 NeRF 表示更直接。

\subsection{用于 SLAM 的 3D Gaussian Splatting}

3DGS 在实时新视图合成方面取得了显著成功，其用于 SLAM 的潜力目前正在积极探索。与 NeRF 类似，3DGS 通过可微渲染的端到端优化，将定位和建图无缝统一。这种集成不再需要多阶段手工设计的算法，并能够建模复杂光场，包括透明物体产生的光场，而这些光场很难用经典三维表示捕获。

\begin{figure}[htbp]
\centering
\includegraphics[width=\linewidth,height=0.78\textheight,keepaspectratio]{6463b73015053972ee817424db46e32750ee327a8b81c868cf047d9bc6cf78e1.jpg}
\caption{图 14.16 Gaussian Splatting 用于语义场景重建的示例。多视图二维图像特征通过可微渲染融合到三维高斯中。图像来自 \textsuperscript{[1316]}（©2024 Springer）。}
\label{fig:chapter-14-6463b73015053972ee817424db46e32750ee327a8b81c868cf047d9bc6cf78e1}
\end{figure}

\begin{figure}[htbp]
\centering
\includegraphics[width=\linewidth,height=0.78\textheight,keepaspectratio]{fee60eac68dbd121ea9be473ac39f7b90cbeb5013b653f065e0ddb0bf7aea1f7.jpg}
\caption{图 14.17 基于 3D Gaussian Splatting 的视觉 SLAM 方法示例 \textsuperscript{[1214, 548, 740]}。图像来自 \textsuperscript{[1214]}（©2024 IEEE）和 \textsuperscript{[548]}（©2024 IEEE）。}
\label{fig:chapter-14-fee60eac68dbd121ea9be473ac39f7b90cbeb5013b653f065e0ddb0bf7aea1f7}
\end{figure}

除这些相似性外，3DGS 相比 NeRF 还具有若干独特优势。第一，其高效的基于光栅化流水线显著加快图像渲染，直接惠及需要实时地图交互的下游应用（例如机器人导航）。第二，其离散、点状的高斯基元表示允许灵活分配和更新场景元素。这种灵活性尤其适合 SLAM：状态估计需要根据运行时观测动态适应，例如更新场景边界和闭合回环。

利用这些属性，3DGS 可以作为视觉 SLAM 系统的核心场景表示（图 14.17）。若干早期工作探索了其同步定位与建图能力 \textsuperscript{[1214, 548, 740]}。一种直接应用是 RGB-D SLAM，其中外部深度传感器（例如飞行时间（ToF）深度相机）提供准确、稠密的深度数据，可作为高斯位置的先验或真实值信号。加入该深度信息后，可以准确初始化三维高斯的位置，并通过最小化深度渲染损失进一步正则化，从而实现准确的相机跟踪以及高保真、快速的地图渲染。

然而，这些方法需要额外硬件才能重建三维高斯，而原始 3DGS 只需要普通相机采集 RGB 图像。理想情况下，方法应完全利用三维高斯灵活且可优化的特性，仅在基于视觉的单目 SLAM 设置中运行。在这一背景下，MonoGS \textsuperscript{[740]} 等工作展示了仅使用彩色图像执行单目 SLAM 的可行性（图 14.18）。在 MonoGS 中，虽然三维高斯的初始位置未知，但可以随机初始化高斯，再通过多视图优化逐步估计和细化其位置。

\begin{figure}[htbp]
\centering
\includegraphics[width=\linewidth,height=0.78\textheight,keepaspectratio]{4df2a894c06e635e1193a391c10fde40bd29d2e80894bbfb6e92cddad90ee25c.jpg}
\label{fig:chapter-14-4df2a894c06e635e1193a391c10fde40bd29d2e80894bbfb6e92cddad90ee25c}
\end{figure}

\begin{figure}[htbp]
\centering
\includegraphics[width=\linewidth,height=0.78\textheight,keepaspectratio]{a118cbdcf8cf1bbdf7c2ac5cdce2da21d1da38015dd62bc09c3de03125c7337b.jpg}
\label{fig:chapter-14-a118cbdcf8cf1bbdf7c2ac5cdce2da21d1da38015dd62bc09c3de03125c7337b}
\end{figure}

\begin{figure}[htbp]
\centering
\includegraphics[width=\linewidth,height=0.78\textheight,keepaspectratio]{d455a6276650a217a212d7505a9600efb1079d88a69af91bae78226c36a06804.jpg}
\caption{图 14.18 MonoGS 的单目重建结果 \textsuperscript{[740]}。得益于三维高斯空间分配的灵活性，该方法无需任何深度监督，仅从彩色图像联合估计相机位姿和场景重建。}
\label{fig:chapter-14-d455a6276650a217a212d7505a9600efb1079d88a69af91bae78226c36a06804}
\end{figure}

三维高斯的离散特性对全局地图更新具有显著优势，尤其是在回环场景中。由于表示由一组独立且定义明确的基元组成，即使环境随时间变化，也更容易匹配和对齐场景的不同片段。例如，LoopSplat \textsuperscript{[1308]} 展示了这一思想的实际应用：离散三维高斯元素充当鲁棒特征，可以在不同时间和视点之间可靠配准。3DGS 作为 SLAM 表示的另一项显著优势，是能够捕获复杂光场。通过有效捕获环境中复杂的光分布，3DGS 不仅提高了视觉保真度，还提供了可改善建图和定位性能的附加线索。I$^2$-SLAM \textsuperscript{[42]} 充分利用这一属性，将运动模糊、色调映射等更精确的基于物理的渲染过程集成到现有 3DGS 框架中，实现 HDR 地图重建与准确相机位姿跟踪的同步进行（图 14.19）。

\begin{figure}[htbp]
\centering
\includegraphics[width=\linewidth,height=0.78\textheight,keepaspectratio]{ab79803e5902c5185edd2b0ce9bcf8f42ea2c75faa857e63b91309e2064d373d.jpg}
\caption{图 14.19 I$^2$-SLAM 的系统概览 \textsuperscript{[42]}（©2024 Springer）。该方法将物理成像过程集成到 SLAM 中，以捕获 HDR 辐射场。}
\label{fig:chapter-14-ab79803e5902c5185edd2b0ce9bcf8f42ea2c75faa857e63b91309e2064d373d}
\end{figure}

尽管 3D Gaussian Splatting（3DGS）作为统一 SLAM 表示的全部潜力仍在积极研究中，另一个有前景的方向是将 3DGS 集成到现有点云处理方法中。只需忽略 3DGS 的附加特征，就可以将其视为点云，因此能够直接结合稳定且在实际应用中性能良好的成熟点算法（图 14.20）。例如，Photo-SLAM \textsuperscript{[488]} 将稀疏跟踪系统（ORB-SLAM）与 3DGS 结合；GS-ICP SLAM \textsuperscript{[424]} 和 RTG-SLAM \textsuperscript{[868]} 则利用 ICP 进行跟踪。在 RTG-SLAM 中，系统进一步加入基于 ORB 关键点的后端优化，取得了快速、准确且具有较高实用性的性能。

此外，3DGS 的固有通用性也延伸到与 LiDAR 测量的融合。近期若干方法提出了 LiDAR+3DGS 混合框架 \textsuperscript{[471, 1054, 624]}，表明 3DGS 可以作为跨多种传感器模态的统一地图表示。从 RGB SLAM 到 LiDAR 增强重建的广泛应用体现了 3DGS 的关键优势：灵活性和适应性使其适用于多样的 SLAM 场景。

\section{延伸阅读与近期趋势}

基于 NeRF/3DGS 的 SLAM 具有端到端优化、复杂光场捕获、实时渲染和动态基元分配等优势，是有前景的研究方向。然而，仍有许多重要技术挑战。本节讨论这些局限，并概述新的趋势和潜在未来研究方向。

\begin{figure}[htbp]
\centering
\includegraphics[width=\linewidth,height=0.78\textheight,keepaspectratio]{a00b530a6134404d08935590671b3d09447e2453850639e2f7a91b36a0256c95.jpg}
\label{fig:chapter-14-a00b530a6134404d08935590671b3d09447e2453850639e2f7a91b36a0256c95}
\end{figure}

(a) PhotoSLAM \textsuperscript{[488]}

\begin{figure}[htbp]
\centering
\includegraphics[width=\linewidth,height=0.78\textheight,keepaspectratio]{81821838f7e5cb4da5184295038472b14ab830364220c4baaf156bea85f55d59.jpg}
\label{fig:chapter-14-81821838f7e5cb4da5184295038472b14ab830364220c4baaf156bea85f55d59}
\end{figure}

(b) RTG-SLAM \textsuperscript{[868]}
图 14.20 将 3DGS 与基于关键点的 SLAM 结合的方法示例。图像来自 \textsuperscript{[488]}（©2024 IEEE）和 \textsuperscript{[868]}（©2024 ACM）。

\textbf{实时性能。} 基于 NeRF/3DGS 的端到端 SLAM 面临的主要挑战之一是处理速度极慢。2015 年，ElasticFusion \textsuperscript{[1184]} 等经典稠密 SLAM 方法已能在 GTX 780Ti 上实现包含跟踪、建图和渲染的实时性能；但大多数基于 NeRF/3DGS 的 SLAM 方法即使在 FLOPs 和显存明显更高的现代 GPU（例如 NVIDIA RTX 4090）上，也难以超过 5 fps。这一限制使当前基于 NeRF/3DGS 的 SLAM 难以用于实时机器人应用，只能依赖离线服务器 GPU。由于 3DGS 支持实时渲染，如何高效优化它仍是关键挑战。提高效率的一条潜在路线是采用紧凑表示，减少高斯数量，从而降低计算开销；另一种方法是使用高阶优化器替代标准一阶梯度下降，以加速收敛。还可以训练网络，通过一次前向传播同时预测相机位姿和三维高斯，从而消除迭代优化。PixelSplat \textsuperscript{[169]} 和 MVSplat \textsuperscript{[191]} 已证明，前馈网络能够从带位姿图像预测三维高斯参数，无需测试时优化。尽管这些早期工作需要外部模块估计的相机位姿，随着 DUSt3R \textsuperscript{[1164]} 的成功（详见第 13 章），不依赖位姿输入、能鲁棒预测点图的无位姿三维重建方法开始受到关注。在此基础上，No-PoSplat \textsuperscript{[1242]}（图 14.21）进一步实现了无位姿 3DGS 预测，展示出有前景的结果。将这些技术用于多视图增量式 SLAM 设置，是值得探索的未来方向。

\textbf{相机跟踪精度。} 基于 NeRF/3DGS 的 SLAM 相机跟踪精度往往落后于经典稀疏 SLAM 方法。这些方法依赖多帧之间的光度一致性，在 Replica \textsuperscript{[1040]} 等具有真实深度和彩色图像的合成数据集上表现良好。然而，在真实场景中，由于传感器噪声、卷帘快门等未建模效应，该假设常常失效。提高精度可能需要更精确地建模场景光场和渲染过程，或探索彩色图像域之外的替代匹配范式。正如上一段所述，引入更鲁棒的前馈预测模型也可能解决这一挑战。

\begin{figure}[htbp]
\centering
\includegraphics[width=\linewidth,height=0.78\textheight,keepaspectratio]{33dc0010dbcad55564bcb2427050ee28aee064a13a81bb7bb2e72f4c9561b773.jpg}
\caption{图 14.21 NoPoSplat 的系统概览 \textsuperscript{[1242]}。该方法直接从两个输入视图预测三维高斯，可用于位姿估计和新视图合成。图像来自 \textsuperscript{[1242]}。}
\label{fig:chapter-14-33dc0010dbcad55564bcb2427050ee28aee064a13a81bb7bb2e72f4c9561b773}
\end{figure}

\textbf{四维场景重建。} 当前大多数 NeRF/3DGS SLAM 方法假设环境静态，但真实场景往往存在显著动态变化。一种常见做法是屏蔽潜在动态物体，只重建静态背景，以获得鲁棒里程计估计。不过，考虑到 3DGS 在非刚性/四维场景重建和点跟踪方面取得的显著成功，探索完整的 4D-SLAM 框架似乎是一条有前景的路线 \textsuperscript{[138, 986, 1289, 741]}。

\textbf{大规模/室外场景重建。} 当前大多数基于 NeRF/3DGS 的 SLAM 方法局限于房间尺度序列，因此大规模和室外场景重建仍是重大挑战。扩展规模需要高效地图更新，例如采用自适应分辨率管理计算负载，并进行漂移校正以维持大范围区域的精度。室外环境还引入额外复杂性，包括天气和一天中不同时间导致的光照变化，这会使光场建模更加困难。
